\documentclass{subfiles}
\begin{document}
    Let's now suppose that need to minimize (resp. maximize) more then a single cost funtion,
        that is, we have an instance of a problem $(X, c_{1}, \ldots, c_{k})$.
        An optimal solution would be an element $\vec{x}^{\mathrm{id}} \in X^{k}$
        such that 
        $$
            x_{i}^{\mathrm{id}} = \min\limits_{x \in X} c_{i}(x)
        $$
        for $i = 0, \ldots, k$. The naive approach,
        would be to define a \emph{fitness} function as linear combination of the cost functions.
        That is, it would the define the function
        $$
            \mathrm{fit}(x) = \sum\limits_{n = 1}^{k} a_{i} c_{i}(x)
        $$
        where $a_{i}$ is the weight of the $i$-th objective.

    The issue with this approach is that we would need a way to determine the weights,
        and this task is non-trivial.

    A better approach consists in the computation of the so-called \emph{Pareto} front.
        Let $\vec{y}_{1}, \vec{y}_{2}$ be two solutions,
        we say that $\vec{y}_{1}$ Pareto-dominates $\vec{y}_{2}$ if $\vec{y}_{1}$ is as good as $\vec{y}_{2}$,
        and there exists an objective for which $\vec{y}_{1}$ is better then $\vec{y}_{2}$.
        In this case we write $\vec{y}_{i} \succ \vec{y}_{2}$.
        In \Cref{Proc:ParetoDomiComputation} we give the pseudo-code to compute the Pareto-dominance.
        \subfile{../TikZ/Procedure - Pareto Dominance}

\end{document}
